%! Summary Statistics for a Quantitative Variable — Unit 1, Lesson 1.5 — Warm-Up
% ONE PAGE, blank and key. Three items, set at 12pt so they fill the page instead of leaving
% the bottom third empty. Do not add a fourth item — the page is full.
% GRADUAL RELEASE: the warm-up activates prior knowledge before the guided notes. Item 1 is
% the spiral back to Lesson 1.4 (name the shape of a pre-drawn dotplot from its tail, and the
% value that stands apart); item 2 averages six numbers and finds the middle one on a set where
% the two AGREE — notes section 1 breaks that agreement on the nine paychecks; item 3 is the
% quartile split done by hand before it has a name — notes section 1 names it.
% NO SKETCH: the item 1 dotplot is pre-drawn; students read it.
\documentclass[12pt]{article}
\usepackage{apstats-article}
\usepackage{apstats-boxes}

\begin{document}

\pageheader{Unit 1, Lesson 1.5}{Warm-Up}

\vspace{0.22in}

\begin{enumerate}[leftmargin=1.9em, itemsep=0.40in]

  \item \textbf{From last lesson.} Twelve students reported how many hours of television they
        watched last week.

        \vspace{4pt}
        \begin{center}
        \begin{tikzpicture}[scale=0.95]
          \draw[->, thick] (-0.2,0) -- (7.1,0);
          \foreach \x/\lab in {0/0, 0.9/2, 1.8/4, 2.7/6, 3.6/8, 4.5/10, 5.4/12, 6.3/14} {
            \draw (\x,-0.07) -- (\x,0.07); \node[below, font=\scriptsize] at (\x,-0.1) {\lab};}
          \foreach \x/\n in {0.45/2, 0.9/3, 1.35/2, 1.8/2, 2.25/1, 2.7/1, 6.3/1} {
            \foreach \k in {1,...,\n} { \fill[navy] (\x,{0.2*\k}) circle (0.075); }}
          \node[font=\scriptsize] at (3.4,-0.55) {Hours of television};
        \end{tikzpicture}
        \end{center}
        \vspace{4pt}

        Describe the \textbf{shape} of this distribution, and name any value that stands apart
        from the rest.\\[14pt]
        \writeline\\[14pt]\writeline

  \item Here are six numbers: \quad 4 \quad 9 \quad 6 \quad 12 \quad 7 \quad 10

        \vspace{12pt}
        Their \textbf{average}: \blank{2.6cm}

        \vspace{12pt}
        Put them in order and give the value in the \textbf{middle}: \blank{2.6cm}

  \item Here are seven numbers, already in order.

        \vspace{8pt}
        \begin{center}
        3 \qquad 5 \qquad 8 \qquad 9 \qquad 12 \qquad 15 \qquad 20
        \end{center}
        \vspace{8pt}

        \textbf{(i)} Which one is in the middle? \blank{2.0cm} \quad
        How many are below it? \blank{1.4cm} \quad Above it? \blank{1.4cm}
        \vspace{12pt}

        \textbf{(ii)} Look only at the numbers \emph{below} the middle one. Which of
        \emph{those} is in the middle? \blank{2.0cm}
        \vspace{12pt}

        \textbf{(iii)} Look only at the numbers \emph{above} the middle one. Which of those
        is in the middle? \blank{2.0cm}

\end{enumerate}

\end{document}
